\chapter[1990 --- Murray et al.]{1990: Joseph E. Murray, E. Donnall Thomas}
\label{ch:1990_Murray}

\section*{Main Contribution}

\textit{Pioneered organ transplantation techniques that made kidney and bone marrow transplantation safe and effective treatments for terminal diseases.}

\section{Official Citation}

for their discoveries concerning organ and cell transplantation in the treatment of human disease

\section{Background and Historical Context}

The story of organ transplantation is one of medicine's most audacious quests---the dream of replacing a failed organ in one human being with a healthy organ from another. By the time the Nobel Committee awarded the 1990 Prize in Physiology or Medicine to Joseph E. Murray and E. Donnall Thomas, this dream had become clinical reality, saving hundreds of thousands of lives worldwide. Yet the road from early experimental surgery to routine transplantation was paved with decades of immunological mystery, surgical ingenuity, and painstaking clinical observation.

Joseph Edward Murray was born on April 1, 1919, in Milford, Massachusetts. He attended the College of the Holy Cross in Worcester and then enrolled at Harvard Medical School, where he developed an early interest in surgery. During World War II, Murray served in the U.S. Army at Valley Forge General Hospital in Pennsylvania, where he treated burn victims and became fascinated by skin grafting---an early encounter with the problems of tissue compatibility that would define his career. After the war, he returned to Boston and completed his surgical training at Brigham Hospital, which would later become part of the Brigham and Women's Hospital complex.

E. Donnall Thomas was born on March 15, 1920, in Mart, Texas. He studied medicine at the University of Texas and then pursued research training at various institutions. Thomas became interested in the problem of bone marrow transplantation after observing patients who developed bone marrow failure following radiation exposure. He moved to the City of Hope National Medical Center in Duarte, California, where he would conduct much of his pioneering work on hematopoietic cell transplantation.

The scientific backdrop against which both men worked was shaped by decades of immunological inquiry. In the early twentieth century, Alexis Carrel and others had demonstrated that organ transplantation between animals was possible but that the transplanted organ was invariably rejected. The immunological basis of rejection was slowly elucidated through the mid-twentieth century: researchers established that the immune system recognized foreign tissue as ``non-self'' and mounted a coordinated attack against it. The major histocompatibility complex (MHC), a set of cell-surface proteins that serve as identity markers, was identified as the primary target of the rejection response. Understanding and circumventing this barrier was the central challenge facing transplantation pioneers.

\section{The Discovery/Contribution}

\subsection{Joseph E. Murray and the First Successful Kidney Transplant}

Murray's contribution began in earnest on December 23, 1954, at Peter Bent Brigham Hospital in Boston, when he performed the first successful human kidney transplant between identical twins---Ronald and Richard Herrick. The operation was a landmark in surgical history. Because the twins were monozygotic and therefore immunologically identical, the transplanted kidney was not rejected. The patient, Richard Herrick, received a kidney from his brother and went on to live another eight years before dying of complications related to the original disease that had destroyed his own kidneys.

The success of the twin transplant validated the concept that organ transplantation could work if the immunological barrier could be overcome. However, Murray immediately recognized that the real challenge lay in transplanting kidneys between non-identical individuals---the vast majority of potential recipients. Through the 1950s and 1960s, Murray conducted a systematic investigation of immunosuppressive strategies. He experimented with total body irradiation, corticosteroids, antimetabolites, and the antilymphocyte serum to suppress rejection responses.

Murray's surgical technique evolved rapidly. He refined the vascular anastomosis---the delicate connection of blood vessels between donor kidney and recipient---and developed methods for preserving the kidney during transfer. His meticulous attention to surgical detail, combined with an aggressive approach to immunosuppression, enabled transplants between non-identical siblings and eventually unrelated donors.

In 1962, Murray performed a kidney transplant using a cadaver organ---a breakthrough that vastly expanded the pool of potential donors. By the mid-1960s, cadaver kidney transplantation was becoming a standard procedure in leading centers around the world, and Murray was at the forefront of establishing the clinical protocols that would govern organ allocation, donor matching, and post-transplant care.

\subsection{E. Donnall Thomas and Bone Marrow Transplantation}

Thomas approached the problem of transplantation from a different angle. Rather than transplanting whole organs, he focused on hematopoietic stem cells---the precursors of all blood cells. The concept was that if a patient's bone marrow was destroyed (by disease, chemotherapy, or radiation), it might be replaced by infusing healthy marrow from a donor.

Thomas's early work in the 1950s was largely unsuccessful. Dogs that received bone marrow transplants after total body irradiation either rejected the graft or developed graft-versus-host disease (GVHD), a devastating condition in which donor immune cells attack the recipient's tissues. Thomas persisted nonetheless, refining his conditioning regimens and developing methods for matching donor and recipient.

His breakthrough came in the late 1960s and early 1970s. Thomas demonstrated that meticulous HLA (human leukocyte antigen) typing of donors and recipients dramatically reduced the incidence and severity of both graft rejection and GVHD. He established that cyclophosphamide and total body irradiation, given before transplantation, could effectively destroy the recipient's immune system and make room for donor stem cells. He showed that treating GVHD with methotrexate could be life-saving.

In 1969, Thomas and his colleague at the Fred Hutchinson Cancer Research Center in Seattle, Rainer Storb, reported a landmark series of bone marrow transplants in patients with leukemia. The results were mixed but demonstrated proof of principle: some patients achieved long-term remission. Thomas continued to refine the procedure throughout the 1970s and 1980s, establishing bone marrow transplantation as the treatment of choice for a range of hematological malignancies and bone marrow failure syndromes.

Thomas also pioneered the concept of allogeneic bone marrow transplantation---using marrow from a matched sibling or unrelated donor---and developed the conditioning regimens, GVHD prophylaxis protocols, and supportive care strategies that remain the foundation of modern transplant practice. His work demonstrated that the transplanted immune system could mediate a graft-versus-leukemia effect, in which donor T cells recognized and destroyed residual cancer cells in the recipient.

\section{Impact on Modern Medicine}

The impact of Murray's and Thomas's work on modern medicine is difficult to overstate. Organ transplantation has become a routine procedure in hospitals around the world. In the United States alone, approximately 40,000 organ transplants are performed each year, with kidney transplantation being the most common. Advances in immunosuppressive drugs---including cyclosporine (introduced in the 1980s), tacrolimus, mycophenolate mofetil, and monoclonal antibodies---have dramatically improved graft survival rates. Today, the one-year survival rate for kidney transplants exceeds 95\%, and many patients live 20 years or more with their transplanted organ.

Bone marrow transplantation, now more accurately called hematopoietic stem cell transplantation (HSCT), has similarly transformed the treatment of hematological diseases. It is the standard of care for many forms of leukemia, lymphoma, aplastic anemia, and inherited blood disorders such as sickle cell disease and thalassemia. The development of reduced-intensity conditioning regimens has made transplantation available to older patients who would not have tolerated the aggressive protocols of earlier decades. The growing use of unrelated volunteer donors, facilitated by international registries such as the National Marrow Donor Program, has expanded access to transplantation for patients who lack matched siblings.

Beyond kidney and bone marrow transplantation, Murray's and Thomas's work laid the groundwork for transplantation of the heart, liver, lung, pancreas, and intestine. Each of these organs required unique surgical techniques and immunological strategies, but all relied on the foundational principles established by the two laureates.

Modern research continues to build on their legacy. Efforts to induce immune tolerance---the ability of the recipient to accept a transplanted organ without ongoing immunosuppression---represent one of the most exciting frontiers in transplantation biology. Xenotransplantation, the transplantation of organs from animals (particularly genetically modified pigs) into humans, is being actively pursued as a potential solution to the organ shortage crisis. The development of artificial organs and bioengineered tissues also draws on the immunological insights that Murray and Thomas helped establish.

\section{Legacy}

Joseph Murray continued to practice surgery at Brigham and Women's Hospital until his retirement. He received numerous honors, including the National Medal of Science in 1990 and the Albert Lasker Clinical Medical Research Award. Murray remained deeply humble about his contributions, often deferring credit to his patients and colleagues. He passed away on November 26, 2012, at the age of 93.

E. Donnall Thomas remained active in research at the Fred Hutchinson Cancer Research Center until his retirement. He received the Albert Lasker Clinical Medical Research Award in 1982, eight years before the Nobel Prize. Thomas continued to advocate for bone marrow transplantation research and for improved access to transplantation for patients worldwide. He died on October 20, 2012, at the age of 92.

Together, Murray and Thomas transformed transplantation from an experimental curiosity into one of a major therapeutic modalities in medicine. Their work embodies the spirit of the Nobel Prize---the translation of fundamental scientific insight into treatments that alleviate human suffering. The millions of patients who have received transplants and gone on to lead healthy, productive lives stand as enduring testimony to their vision, courage, and perseverance.


\section{Modern Developments}

Murray and Thomas' transplant work has evolved into modern transplantation medicine. Tacrolimus and mycophenolate have improved immunosuppression. Belatacept provides costimulation blockade. Understanding of regulatory T cells has led to tolerance-inducing protocols. Desensitization protocols enable ABO-incompatible transplantation. Machine perfusion extends organ viability. Understanding of transplant immunology has informed development of tolerance-inducing therapies that may eliminate the need for lifelong immunosuppression.
